\section{问题重述与分析}
\subsection{研究对象与任务}
题目研究由光伏、负载、储能和外部电网构成的非孤岛式微网\cite{problem}。光伏供给与负载需求在时间上并不一致，电价又随时段变化。储能能够转移电量，但受到能量容量、充放电功率和转换损耗限制。购电不足时的紧急补给以及日内修改购电计划产生的费用，使“按预计缺口购电”不能直接等同于总费用最小的策略。

依据四问主模型及其正式结果\cite{models}，问题一讨论已知典型日负载、光伏与电价时的确定性调度；问题二在实际负载和光伏尚未实现时制定日前计划；问题三进一步利用0、6、12、18点发布的光伏预报调整决策；问题四分别在第二、三问框架中引入波动电价。本文保留五组既有主结果，不以其他历史候选替换主方案。

\subsection{建模思路与各问衔接}
第一问用于建立共同物理模型并解释储能收益。第二问的关键是平衡常规购电与高价紧急购电：过多计划购电会产生闲置费用，过少则增加缺口风险，因此以历史误差场景近似缺口的期望损失。跨日前瞻用于保留日末电量对次日的价值，避免人为逐日清空或复位电池。

第三问新增的信息并非实际未来光伏，而是逐次发布的预报。模型需要同时回答如何融合预报、如何表达未来信息分支，以及何时为修改计划支付费用。第四问的新增难点是电价与缺口可能相关，应在同一结算口径下分析费用变化。4-2继承日前决策，4-3继承多时点调整；二者的初始化与决策权限不同，不能仅凭总费用判断哪一算法更优。

\section{统一假设、数据口径与符号}
\subsection{共同约定与局部条件}
本文按以下约定建立主模型。
\begin{enumerate}
\item 每天划分为144个十分钟时段，段内功率视为平均功率。时间标签采用区间末时刻，00:10表示00:00—00:10，电量由功率乘$\Delta t=1/6$小时得到。
\item 不向外网售电；允许弃光。随机调度中允许弃置多余供给，已计划但未使用的购电仍计费。未加入折旧或寿命成本。
\item 题给额定容量12000 kWh、运行储电量1200—10800 kWh及功率上限5000 kW。主模型将功率限制解释为设备外端限制，并将“充放电效率90\%”解释为充电与放电效率各0.9，即往返0.81。另一效率口径在第\ref{sec:efficiency}节单独讨论。
\item 日前决策仅使用决策日前已实现的数据；日内决策可使用已发布预报和截至当时可见的信息。未来实际值仅用于事后结算，不用于选择当前控制。
\item 紧急购电按对应供电时段实际电价的5倍结算。第三问和4-3的取消量退回原费、另付50\%违约费，增购按1.5倍计价。该取消规则是本模型明确采用的结算解释。
\end{enumerate}

附件1对应典型日电价、负载和光伏预测曲线，第一问将该预测作为确定性输入；附件2对应全年实际负载和光伏，附件3对应发布预报，附件4对应实际波动电价。附件5规定五份结果文件及工作表内容。这里沿用主模型已有的时间转换、历史预测和样本处理，不额外假定完成了未记录的数据清洗。

\begin{table}[htbp]\centering\small
\caption{五组结果的统计期与储电量条件}\label{tab:conditions}
\begin{tabularx}{\linewidth}{l X r r}
\toprule 结果集 & 统计期及初始化 & 起始电量/kWh & 终止电量/kWh\\\midrule
问题一 & 典型日，日初固定 & 6000 & 6000\\
问题二、4-2 & 2—12月；1月共同调度预热 & 8550 & 6000\\
问题三、4-3 & 2—12月；1月保持电量并积累历史 & 6000 & 6000\\
\bottomrule\end{tabularx}\end{table}

附件5要求年度输出覆盖2025年2月1日至12月31日，共334天，1月用于主模型预热。题给2025年1月1日0点储电量6000 kWh；第一问另明确日初、日末相等。表\ref{tab:conditions}中的2月初状态及年度模型的年末6000 kWh约束是现有模型的初始化和终端约定，不冒充题面逐字要求。第三问使用统一6000 kWh起点的历史预测对照，而非直接把第二问正式费用作为其同条件基准。一般滚动窗口末端条件也不同于每天日末回到日初。

题面还要求在论文中按表1、表2、表3给出指定结果，并讨论是否需要其他时刻的预报。本文在各问呈现总费用与图示，在附录\ref{sec:outputs}完整列示指定十分钟购电、四小时充放电和连续紧急购电区间；日期为3月20日、6月21日、9月23日及12月21日。第三问的更新频率比较直接回应预报时刻选择。第四问按题意重算第二、三问，不额外设立4-1。

\subsection{符号与储能约束}
\begin{table}[htbp]\centering\small
\caption{主要符号与单位}
\begin{tabularx}{\linewidth}{l X l}\toprule
符号 & 含义 & 单位\\\midrule
$d,t,s,n$ & 日期、日内时段、历史场景、场景树节点索引 & —\\
$L_t,V_t,N_t$ & 负载、光伏及净负载电量，$N_t=L_t-V_t$ & kWh\\
$\widehat N_t,N_{s,t}$ & 净负载点预测与场景值 & kWh\\
$g_t,g_t^0,q_t$ & 日前购电、0点原计划、最终调整后购电 & kWh\\
$c_t,b_t,E_t$ & 外端充电量、外端放电量、段末内部储电量 & kWh\\
$e_t,w_t$ & 紧急购电量、弃光或弃置供给量 & kWh\\
$u_t,v_t$ & 相对原计划的增购量与取消量 & kWh\\
$p_t,\widehat p_t,p_{s,t}$ & 实际、预测与场景电价 & 元/kWh\\
$\omega_s,\pi_n$ & 场景权重与节点概率 & —\\
$\eta_c,\eta_d$ & 充电效率、放电效率 & —\\
\bottomrule\end{tabularx}\end{table}

为避免日期索引$d$与放电量混淆，全文用$b_t$表示主程序中的放电变量。共同储能约束写为
\begin{align}
E_t&=E_{t-1}+\eta_c c_t-b_t/\eta_d,\label{eq:storage}\\
1200\le E_t&\le10800,\qquad 0\le c_t,b_t\le5000/6.\label{eq:bounds}
\end{align}
式\eqref{eq:storage}区分设备外端电量与内部库存，不能直接把充电量减放电量作为库存变化。跨日时下一天的起点必须等于上一天实际执行后的末电量。
